% 1. Forzamos números arábigos y empezamos sin encabezado (para el índice)
\pagenumbering{arabic}
\pagestyle{empty}

% 2. Definición del código del documento como comando
\newcommand{\docCode}{SIGCHOS-MNL-V1-2026-001}

% 3. Definición del comando de activación
\newcommand{\activarEncabezado}{
    \pagestyle{fancy}
    \fancyhf{}

    % --- ENCABEZADO: tabla de 3 columnas con borde negro dividiendo ---
    \fancyhead[C]{%
        \renewcommand{\arraystretch}{1.2}%
        \setlength{\tabcolsep}{4pt}%
        \begin{tabularx}{\textwidth}{|m{2.5cm}|>{\centering\arraybackslash}X|m{5.5cm}|}
            \hline
            \multirow{3}{*}{\centering\includegraphics[width=2.5cm,keepaspectratio]{images/logo-espe.png}} & 
            \multirow{3}{*}{\textbf{\fontsize{10pt}{12pt}\selectfont\makecell{MANUAL TÉCNICO DEL\\[2pt]SISTEMA DE COMERCIO ELECTRÓNICO}}} & 
            \textbf{\fontsize{9pt}{10pt}\selectfont ÁREA DE TECNOLOGÍA E INNOVACIÓN} \\ \cline{3-3}
            & & \fontsize{8pt}{9pt}\selectfont Código de documento:\ \docCode \\ \cline{3-3}
            & & \fontsize{9pt}{10pt}\selectfont Página:\ \thepage\ de \pageref*{LastPage} \\
            \hline
        \end{tabularx}%
    }

    % Pie de página: código de proceso (izq) y revisión (der)
    \fancyfoot[L]{\footnotesize\textit{Código de proceso: SIG.1.0}}
    \fancyfoot[R]{\footnotesize\textit{Rev. ATI: 2026-abr-23}}


    % Aplicar a páginas de inicio de capítulo
    \fancypagestyle{plain}{\pagestyle{fancy}}
    \renewcommand{\headrulewidth}{0pt}
}